\section{Заключение}

В ходе выполнения лабораторной работы изучены основные возможности
системы контроля версий Git и приобретены практические навыки работы с ней
через графический интерфейс Git GUI/gitk. 

В рамках работы был создан локальный репозиторий, продемонстрирована
фиксация изменений (в том числе выборочная – с отменой отдельных строк
перед коммитом), создание и переключение веток, слияние веток с разрешением
конфликта, просмотр истории изменений и авторства строк кода, а также отмена
изменений через revert. Репозиторий опубликован на удалённом сервере GitHub,
продемонстрирована синхронизация локальных и удалённых изменений.

В рамках дополнительного задания освоены создание и применение тегов,
отмена изменений командами revert и reset, перебазирование веток с
разрешением конфликтов, временное сохранение и восстановление изменений
с помощью stash, а также создание и применение серии патчей.

Полученные навыки позволяют организовать коллективную разработку
программного проекта с сохранением полной истории изменений и
возможностью параллельной работы над разными частями функционала.
